\documentclass[11pt]{amsart}
\usepackage[T1]{fontenc}
\usepackage{lmodern}
\usepackage{amsmath,amssymb}
\usepackage[margin=1in]{geometry}
\usepackage[colorlinks=true,linkcolor=blue,citecolor=blue,urlcolor=blue]{hyperref}
\newcommand{\TC}{\operatorname{TC}}
\newcommand{\rank}{\operatorname{rank}}
\newcommand{\F}{\mathbb F_2}
\newcommand{\R}{\mathbb R}
\newcommand{\Z}{\mathbb Z}
\newcommand{\Q}{\mathbb Q}
\newtheorem{theorem}{Theorem}
\newtheorem{proposition}[theorem]{Proposition}
\newtheorem{lemma}[theorem]{Lemma}
\title[An oriented flat counterexample]{An oriented flat counterexample to the centre-rank formula for topological complexity}
\author{Kaiwen Jin}
\email{yui.ui.siki@gmail.com}
\date{October 13, 2026}
\begin{document}
\begin{abstract}
We construct a closed connected oriented flat five-manifold whose fundamental
group has centre of rank two and whose unreduced topological complexity is at
least ten. Since its universal cover is Euclidean space, this gives a negative
answer to the proposed formula $\TC(M)=2\dim M+1-\rank Z(\pi_1M)$ for closed
oriented aspherical manifolds. The lower bound follows from a nonzero product
of nine degree-one zero divisors in the mod-$2$ cohomology ring of a real Bott
manifold.
\end{abstract}
\maketitle

\section{The centre-rank question}
For a path-connected space $X$, its unreduced topological complexity $\TC(X)$
is the least number of open subsets covering $X\times X$ on which the path
endpoint map $PX\to X\times X$ admits continuous sections. Thus $\TC(X)$
counts the continuous local rules needed for motion planning, with both the
initial and final configuration allowed to vary \cite{Farber}.

For a closed connected oriented aspherical $n$-manifold $M$, put
\[
 z=\dim_{\Q}\bigl(Z(\pi_1M)\otimes_{\Z}\Q\bigr).
\]
The centre-rank question asks whether
\begin{equation}\label{eq:formula}
 \TC(M)=2n+1-z.
\end{equation}
This formula is posed in \cite[Question~5.8]{Neofytidis}. We show that it fails
already among oriented flat manifolds.

\begin{theorem}\label{thm:main}
There is a closed connected oriented flat five-manifold $M$ with contractible
universal cover $\R^5$ such that
\[
 \dim_{\Q}\bigl(Z(\pi_1M)\otimes_{\Z}\Q\bigr)=2,
 \qquad \TC(M)\geq 10.
\]
In particular, $M$ is an aspherical counterexample to \eqref{eq:formula},
whose right-hand side is $9$ for this manifold.
\end{theorem}

The construction uses the affine realization and cohomology presentation of
real Bott manifolds due to Kamishima and Masuda \cite{KM}. A mixed product
of nine zero divisors then gives the lower bound through Farber's theorem
\cite[Theorem~7]{Farber}. An exact value of $\TC(M)$ is not needed for the
counterexample. Section~2 constructs the manifold and computes its centre;
Section~3 gives the cup product and completes the proof. All cohomology below
has coefficients in $\F$, and $\TC$ always uses the unreduced convention.

\section{The manifold and its centre}
Consider the strictly upper triangular binary matrix
\begin{equation}\label{eq:matrix}
 A=\begin{pmatrix}
 0&1&0&1&0\\
 0&0&0&1&1\\
 0&0&0&1&1\\
 0&0&0&0&0\\
 0&0&0&0&0
 \end{pmatrix}.
\end{equation}
For $1\leq i\leq5$, let $e_i$ be the $i$th coordinate vector and define
the affine isometry
\begin{equation}\label{eq:generators}
 s_i(u)=D_i u+\tfrac12 e_i,
 \qquad
 D_i=\operatorname{diag}\bigl((-1)^{A_{i1}},\ldots,(-1)^{A_{i5}}\bigr).
\end{equation}
Set $\Gamma=\langle s_1,\ldots,s_5\rangle$ and $M=\R^5/\Gamma$.
These are exactly the generators in \cite[equation~(3.1)]{KM}, and
\cite[Lemma~3.1]{KM} identifies this quotient with the real Bott manifold
$M(A)$ and its fundamental group with $\Gamma$.

\begin{proposition}\label{prop:geometry}
The manifold $M$ is closed, connected, oriented and aspherical, and
$\dim_{\Q}(Z(\Gamma)\otimes_{\Z}\Q)=2$.
\end{proposition}
\begin{proof}
Since $A_{ii}=0$, the square $s_i^2$ is the unit translation by $e_i$.
The integer translation lattice $\Z^5$ is therefore a normal subgroup of
$\Gamma$. On $T^5=\R^5/\Z^5$ the five generators commute and have order two:
a sign change and a half translation commute modulo integer translations.
For a nonempty product of distinct generators, choose its smallest index $i$.
Every selected generator of larger index fixes coordinate $i$, whereas $s_i$
translates that coordinate by $1/2$. Hence the product has no fixed point.
The induced action is consequently a free action of $(\Z/2)^5$.

It follows that $M=T^5/(\Z/2)^5$ is a smooth compact connected manifold
without boundary. Its Euclidean metric descends to a flat metric. The lifted
action on $\R^5$ is free and properly discontinuous: all translation vectors
lie in $\tfrac12\Z^5$, the linear parts are diagonal sign matrices, and only
finitely many such affine maps can move a compact set to meet itself. Thus
$\R^5$ is its contractible universal cover. Every row of $A$ has even sum,
so every $D_i$ has determinant $+1$ and $M$ is oriented.

Let $g(u)=Bu+v$ be central in $\Gamma$. Commuting with the five unit
translations gives $Be_i=e_i$ for every $i$, so $B=I$. A translation by $v$
commutes with $s_i$ precisely when $D_i v=v$. For the matrix
\eqref{eq:matrix} these conditions are
\[
 v_2=v_4=v_5=0.
\]
The unit translations in coordinates $1$ and $3$ are central, while every
central vector lies in $\tfrac12\Z e_1\oplus\tfrac12\Z e_3$. The two unit
translations therefore span the centre after tensoring with $\Q$ and remain
linearly independent under the coordinate map to $\Q^2$. Its rank is exactly
two.
\end{proof}

\section{A product of nine zero divisors}
By \cite[Lemma~2.1]{KM}, the cohomology ring $R=H^*(M;\F)$ has degree-one
generators $a,b,c,d,e$, corresponding to $x_1,\ldots,x_5$, and presentation
\begin{equation}\label{eq:ring}
 R=\F[a,b,c,d,e]\big/
 \bigl(a^2,\ b^2+ab,\ c^2,\ d^2+(a+b+c)d,\ e^2+(b+c)e\bigr).
\end{equation}
Each successive relation is monic quadratic in the newly adjoined variable.
Polynomial division therefore gives a free rank-two extension with basis
$1,x_j$ at each stage. The squarefree monomials are a basis of $R$; in
particular,
\[
 \alpha=abcd\in H^4(M;\F),\qquad
 \omega=abcde\in H^5(M;\F)
\]
are nonzero.

For $x\in R$, write $\bar x=x\otimes1+1\otimes x$ in the graded algebra
$R\otimes_{\F}R$. This lies in the kernel of cup multiplication
$R\otimes R\to R$, which is the zero-divisors ideal used in
\cite[Definition~6]{Farber}. All graded tensor signs equal $1$ over $\F$.

\begin{lemma}\label{lem:product}
In $R\otimes_{\F}R$,
\begin{equation}\label{eq:product}
 \bar a\,\bar b\,\bar d^{\,5}\,\bar e^{\,2}
 =\alpha\otimes\omega+\omega\otimes\alpha\neq0.
\end{equation}
\end{lemma}
\begin{proof}
Put $u=a+b+c$. The relations \eqref{eq:ring} give
\[
 u^2=ab,\qquad u^3=abc,\qquad u^4=0.
\]
Indeed, $u^3=ab(a+b+c)=abc$, since $a^2=0$ and $ab^2=a^2b=0$.
It follows from $d^2=ud$ that
\[
 d^4=u^3d=\alpha,\qquad d^5=u^4d=0.
\]
The binomial coefficients $\binom5j$ are odd only for $j=0,1,4,5$.
Consequently
\[
 \bar d^{\,5}=\alpha\otimes d+d\otimes\alpha.
\]
Since $a\alpha=b\alpha=0$, multiplication by $\bar a\bar b$ leaves
\begin{equation}\label{eq:intermediate}
 \bar a\bar b\bar d^{\,5}
 =\alpha\otimes abd+abd\otimes\alpha.
\end{equation}
Finally,
\[
 \bar e^{\,2}=e^2\otimes1+1\otimes e^2,
 \qquad \alpha e^2=0,
 \qquad abd\,e^2=abd(b+c)e=\omega.
\]
Multiplying \eqref{eq:intermediate} by $\bar e^{\,2}$ proves the identity.
The two summands have distinct bidegrees $(4,5)$ and $(5,4)$ and are both
nonzero, so their sum is nonzero.
\end{proof}

\begin{proof}[Proof of Theorem~\ref{thm:main}]
Proposition~\ref{prop:geometry} supplies the required manifold and its centre
rank. The product \eqref{eq:product} contains $1+1+5+2=9$ zero-divisor
factors. Farber's lower bound \cite[Theorem~7]{Farber} states, in the
unreduced convention, that $\TC(M)$ is strictly greater than the
zero-divisors cup length of $H^*(M;\F)$. Hence $\TC(M)\geq10$, while
$2\dim M+1-\dim_{\Q}(Z(\pi_1M)\otimes_{\Z}\Q)=9$.
\end{proof}

\section*{AI use}
OpenAI Codex was used for the finite construction search, algebraic checks,
and preparation of the text. The proof is given in full through the displayed
algebraic calculations and the cited cohomology and topological complexity results.

\begin{thebibliography}{9}
\bibitem{Farber}
M.~Farber, \emph{Topological complexity of motion planning},
Discrete Comput. Geom. \textbf{29} (2003), 211--221.
\href{https://doi.org/10.1007/s00454-002-0760-9}{doi:10.1007/s00454-002-0760-9}.
Preprint: \href{https://arxiv.org/abs/math/0111197}{arXiv:math/0111197v1}.
\bibitem{KM}
Y.~Kamishima and M.~Masuda, \emph{Cohomological rigidity of real Bott manifolds},
Algebr. Geom. Topol. \textbf{9} (2009), 2479--2502.
\href{https://arxiv.org/abs/0807.4263v1}{arXiv:0807.4263v1}.
\bibitem{Neofytidis}
C.~Neofytidis, \emph{Topological complexity, asphericity and connected sums},
J. Appl. Comput. Topol. \textbf{9} (2025), article~10.
\href{https://doi.org/10.1007/s41468-025-00205-z}{doi:10.1007/s41468-025-00205-z}.
Preprint: \href{https://arxiv.org/abs/2212.08962v2}{arXiv:2212.08962v2}.
\end{thebibliography}
\end{document}
